%======================================================================
% Section 2: Literature Landscape and Scholarly Context
%======================================================================
\section{Literature Landscape and Scholarly Context}
\label{sec:literature}

\subsection{The organising problem and the shape of the field}

The literature of generalized continuum wave mechanics addresses the infinite propagation speed associated with Fourier heat conduction. Joseph and Preziosi’s retrospective traces early finite-speed heat-flux formulations to Cattaneo and describes Vernotte’s later presentation of the relaxation law as a simple generalization \cite{Joseph1989}.

That pattern --- identification, modification, cost --- is the organising theme of this review, and it is why the literature is best read by theory family rather than by paper. Each family below answers the same question in a different way, and each pays for the answer in a different currency: a second temperature, a second wave, a loss of ellipticity, a loss of uniqueness, a loss of local well-posedness. Sections~\ref{sec:thermoelastic} through~\ref{sec:advanced} take the families in turn; this section establishes how they arose, where the field's centre of gravity has moved, and what the geographic distribution of the work reveals about the structure of the problem.

\subsection{Four phases in the development of the field}

Cattaneo and Vernotte proposed relaxation-type modifications of Fourier heat conduction; Lord and Shulman incorporated a relaxation time into coupled thermoelasticity. Green and Lindsay introduced a temperature-rate-dependent, single-temperature formulation with two relaxation parameters. These models have different constitutive structures and limiting assumptions.

\paragraph{Phase II --- second sound and thermal-displacement theories (1966--1993).} Phonon kinetic theory supplied a microscopic account of non-Fourier transport. Guyer and Krumhansl (1966) derived the phonon-Boltzmann structure and a nonlocal flux length scale; this is not treated as experimental evidence or a continuum thermoelastic law \cite{GuyerKrumhansl1966}. A separate companion paper, “Thermal Conductivity, Second Sound, and Phonon Hydrodynamic Phenomena in Nonmetallic Crystals,” begins at p. 778; it is distinct from the p. 766 article cited for the kinetic-theory equation above \cite{GuyerKrumhansl1966II}. The Green--Naghdi type-II primary formulation is the 1993 nondissipative theory, while the type-III equations are in the 1992 *Journal of Thermal Stresses* paper \cite{Green1993,GreenGN1992}. The type labels and dissipation statements below are limited to those specified systems; no universal Hamiltonian or energy claim is inferred. The synthesis of this phase is the review literature: Chandrasekharaiah's articles on hyperbolic thermoelasticity \cite{Chandrasekharaiah1986,Chandrasekharaiah1998}, Joseph and Preziosi's account of heat waves from kinetic theory to continuum mechanics \cite{Joseph1989,Joseph1990}, and the monographs by Nowacki \cite{Nowacki1975} and Achenbach \cite{Achenbach1973}.

\paragraph{Phase III --- the microscale turn (1995--2011).} Phase II assumed that the Fourier paradox was a matter of taking a limit too early. Phase III questioned the premise, and asked instead what a thermal wave should look like when the material has microstructure. Tzou's dual-phase-lag model \cite{Tzou1995} and Roy Choudhuri's triple-phase-lag formulation \cite{RoyChoudhuri2007} distinguish the flux and temperature-gradient lags and include a thermal-displacement-gradient contribution in the latter; the exact truncation and admissibility conditions remain model-specific. In parallel, Eringen's nonlocal elasticity \cite{Eringen1972,Eringen1974,Eringen1983} and the strain-gradient formulations \cite{Askes2006,Askes2011,Polizzotto2001} transferred the same logic to the mechanical side, and the size-dependent wave speeds measured and computed for nanostructures \cite{Maranganti2006,Maranganti2007} supplied the corresponding experimental evidence.

\paragraph{Phase IV --- memory, fractional operators and multiphysics (2011--present).} Recent models use memory kernels and coupled fields. Gorenflo and Mainardi provide mathematical background on fractional calculus \cite{Gorenflo1997}. Memory-dependent derivatives use time-history weights \cite{Wang2011}; fractional models replace integer-order derivatives with fractional orders in the governing equations \cite{Caputo1967,Povstenko2004,Povstenko2007,Povstenko2009,Povstenko2011,Povstenko2016}. Couplings to electrical, porous, micro-rotational, and carrier fields produce multiphysics formulations \cite{Iesan1986,Iesan2013,Cowin1983,Nunziato1979,Chandrasekharaiah1996,Marotti2013,Sharma2022,Abouelregal2023}.

\subsection{International structure of the field}

The development of generalized thermoelasticity involved several research communities. Foundational heat-conduction and thermoelasticity works, micropolar extensions, anisotropic models, fractional and memory-dependent formulations, and more recent multiphysics studies are represented in the cited literature \cite{Joseph1989,Biot1956,Lord1967,Green1972,Eringen1966a,Eringen1968b,Eringen1972,Dhaliwal1980,Sharma2022,Povstenko2004,Povstenko2007,Povstenko2009,Povstenko2011,Povstenko2016,Wang2011,Iesan1986,Iesan2013,Cowin1983,Nunziato1979,Chandrasekharaiah1996,Marotti2013,Abouelregal2023,Quintanilla2019,Barak2023,Barak2023IJP,Barak2023JTS,Barak2023MMMS,Gupta2023Energy,Gupta2023IJNMF,Kumar2023PhysScr,Gupta2026HeatMap}.

This distribution is not incidental. Each of the centres identified above entered the field through a different scientific interest --- heat conduction in metals, micropolar fluid mechanics, anomalous diffusion in geophysics, piezoelectricity in smart materials --- and each imported the constitutive modification appropriate to that interest. The result is a literature in which the same physical question, the admissibility of a propagating thermal mode, is attacked with structurally different devices: a relaxation time, a second temperature, a thermal displacement, a phase lag, a memory kernel, a fractional order. A review organised around papers would obscure this, because the same device reappears in every centre with different notation. A review organised around the devices --- which is what this review attempts --- shows that the field is not a sequence of independent advances but a single continuing negotiation between three requirements that are individually necessary and jointly expensive: finite propagation speed, thermodynamic admissibility, and well-posedness.

\subsection{The literature corpus}

The studies discussed here range from foundational analyses and surveys to applied work on coupled, nonlocal, fractional, and memory-dependent models. Comparisons are organized by constitutive structure and stated assumptions rather than by bibliographic metadata alone.

\subsection{Development axes and the gaps they expose}

The corpus can be classified along five axes --- physical domain, mathematical structure, modal content, coupling topology and validation regime --- and this classification is developed formally in Section~\ref{sec:taxonomy}. Its value at this stage is diagnostic: it shows where the field is dense and where it is thin. The dense regions are the single-field, one-dimensional, interface-focused problems of classical, relaxation-time and phase-lag thermoelasticity, for which a large body of modal and energy-ratio results exists. The thin regions are more revealing. There is very little work that closes the loop from constitutive novelty to experimental discrimination: the nonlocal and fractional families are theoretically mature, with well-established dispersion relations and modal structures, yet the number of studies that fit predictions to measured dispersion data remains small. The comparison of competing mechanisms for the same physical effect --- nonlocal versus gradient, exponential versus power-law memory, phase lag versus relaxation time --- is similarly under-represented, even though it is precisely in this comparison that the theory choices are testable. A third thin region is the treatment of interfaces in the presence of memory or nonlocality, where the constitutive history and the geometric matching condition must be reconciled, and where the energy partition is not determined by the local constitutive law.

These gaps are not defects of the individual contributions; they follow from the structure of the field, in which each advance has been made by extending a constitutive law and then solving the resulting modal problem, and the natural next step --- putting two such extensions in competition against an experiment --- has not yet become the standard mode of work. Section~\ref{sec:future} takes this diagnosis as the basis for its research directions.

\subsection{Existing reviews and the contribution of this article}

Reviews already exist for several individual theory families: nonlocal elasticity and its application to nanostructures \cite{Askes2006,Askes2011,Maranganti2007}, fractional thermoelasticity \cite{Chandrasekharaiah1998,Povstenko2016}, hyperbolic thermoelasticity and second sound \cite{Chandrasekharaiah1996,Joseph1989}, and generalized thermoelasticity with modified models \cite{Shakeriaski2021}. Their common structure is the one this review sets aside: an enumeration of theory families, each with its governing equations, its application domains and a list of its principal results, with the families treated in isolation and the transitions between them left implicit.

The present review differs in three respects. It is organised by the problems --- finite propagation speed, thermodynamic admissibility, well-posedness, scale dependence, interface partition of energy --- rather than by the theories, so that a reader interested in one problem can follow it across all the theories that address it. It states the cost of each modification alongside the benefit, which the standard format does not do, and which the literature itself rarely does: the second temperature, the loss of ellipticity, the loss of uniqueness of the initial-value problem, and the non-locality of the energy flux are each the price of something that was bought. It separates a mechanism's plausibility from the empirical support for it, treating the evidence base as an explicit axis of classification rather than a matter of emphasis.

\begin{figure}[htbp]
\centering
\includegraphics[width=0.95\textwidth]{fig03-literature-landscape.pdf}
\caption{Selected theory families and relationships in generalized continuum mechanics. This schematic is a conceptual taxonomy of the model families discussed in the text.}
\label{fig:literature_landscape}
\end{figure}

\subsection{Citation and assumptions}

A useful comparison states the validity domain and conventions of each equation and distinguishes a derivation from an attributed result. The assumptions used in each cited formulation are therefore identified alongside the relevant discussion.

Where material errors are conclusively established in the literature, they are recorded transparently and without personal characterisation. Where uncertainty exists, it is acknowledged as an open question. No error allegations are made against any publication, and model-specific differences are never presented as mathematical contradictions without sufficient evidence.

% LocalWords:
